% Partie : lambda-calcul
% Chapitre : lambda-définissabilité
% Section : listes

\documentclass[../../../include/open-logic-section]{subfiles}

\begin{document}

\olfileid[fr]{lam}{ldf}{lst}
\olsection{Listes}

Soient $s$ et $z$ deux variables distinctes qui n'apparaissent libres dans aucun des $M_i$. On peut alors définir la liste $[M_1,M_2,\ldots,M_n]$ par :
\[
  [M_1, M_2, \ldots, M_n] \ident \lambd[sz][s M_1 (s M_2 (\ldots (s M_n z)))]
\]
Informellement, le terme représentant une liste est sa
fonction d'accumulation. Il reçoit deux termes : le premier
est une fonction qui reçoit un élément et une valeur
accumulée, puis renvoie la nouvelle valeur accumulée ;
le second est la valeur accumulée initiale. La fonction
d'accumulation combine les éléments un à un et renvoie
finalement le résultat.

\begin{digress}
Si vous connaissez la programmation fonctionnelle, vous
reconnaîtrez une définition semblable au \emph{pliage à
droite}, ou \texttt{fold\_right} : une liste est définie
comme la fonction de pliage à droite de ses éléments.
\end{digress}

Ce codage permet de définir facilement quelques fonctions
utiles :
\begin{align*}
  \fn{Sum} & \ident
  \lambd[l][l \, (\lambd[xa][\fn{Add} \, x \, a])\,  \num{0}]\\
  \fn{Len} & \ident
  \lambd[l][l \, (\lambd[xa][\fn{Add} \, a \, \num{1}]) \num{0}]
\end{align*}

$\fn{Sum}$ calcule la somme d'une liste de numéraux de Church.
Elle procède par accumulation sur la liste : la valeur
initiale est $\num{0}$ et, pour chaque élément, la nouvelle
valeur accumulée est $\fn{Add}\,x\,a$, où $x$ est l'élément
considéré et $a$ la valeur accumulée. Le résultat est la
somme de tous les éléments.

\begin{prob}
Que fait $\fn{Len}$ ? Expliquer.
\end{prob}

\begin{prob}
Définir une fonction qui renverse une liste.
\end{prob}

\end{document}
